\scanpage{44}
\mspage{18}

b$'$) Cas $\mathscr X\simeq\mathbb P^1_{\mathbb C}$, $\sigma$ une inversion -- laquelle définit une structure de droite projective réelle -- on trouve \deleted{le disque} pour $X$ un \uline{disque}.

Pour obtenir une métrique distinguée sur $\mathscr X$ invariante par $\sigma$, il faut se donner un pt du disque ouvert $X\setminus\partial X$

\hfill\rule{0.34\linewidth}{0.4pt}\hfill\null

En résumé, tte surface conforme à bord $0$-connexe admet une unique métrique \deleted{\unsure{d}}Riemannienne distinguée, sauf dans les quatre cas suivants

a) a$'$) Cas $X\simeq\mathbb C$, ou $X\simeq\deleted{\unsure{\mathbb C}}\mathbb H=\{z\in\mathbb C\mathbin{/}\unsure{J|I|\mathfrak I}z\geq 0\}$ [droite affine complexe, demi plan à bord défini par une droite affine réelle], on a alors une classe de métriques mod constantes, qui forment une orbite sous $\mathrm{Aut}\, X$

b) b$'$) $X\simeq\mathbb P^1_{\mathbb C}$ ou $X\simeq\text{disque}$ [droite projective complexe ou \deleted{\unsure{demi espace projectif|demi plan projectif|droite projective}}\inserted{disque conforme} défini par droite projective réelle] -- il y a alors un ens. de métriques distinguées dépendant de 3 paramètres réels resp de 2 paramètres réels, également orbite sous $\mathrm{Aut}\,X$.

% archive sheet number: 46
